\part{Final evaluation: closing the research objective}\label{part:final}

\begin{keybox}{What this part is}
The last round of work audited the real runtime path, froze and reproduced the previous evaluation, fixed the gaps
it found, and then ran one \textbf{predeclared} evaluation (\code{docs/thesis/final_evaluation_plan.md}):
\begin{itemize}
\item fresh seeds 20--24;
\item \NFRuns{} runs on the final code plus \NOvsRuns{} runs on real Open vSwitch;
\item paired tests with 95\,\% CIs, rank-biserial effect sizes and Holm correction over the primary claims.
\end{itemize}
Everything below comes from that evaluation's artefacts. It supersedes Part~\ref{sec:qoseval}'s v4 numbers, which
are preserved and were reproduced before any change (Q802).
\end{keybox}

\section{Audit and baseline}

\begin{qa}{801}{What did tracing the real runtime path find?}
The audit (\code{docs/RUNTIME_PATH.md}) traced each stage, Monitor $\rightarrow$ Telemetry $\rightarrow$ Classification
$\rightarrow$ Context $\rightarrow$ Demand/prediction $\rightarrow$ Allocation $\rightarrow$ Safety $\rightarrow$
Actuation $\rightarrow$ Verification $\rightarrow$ Learning, to its call site. It found eight gaps:
\begin{itemize}
\item SliceNet decisions did not reach uplink demand estimation;
\item monitoring output fed no decision;
\item there was no validator on the final allocation before actuation;
\item ``feasible'' ignored burst requirements;
\item prediction had never been tested against ramps or an upper bound;
\item there was no truly context-blind control, and controllers were given the workload's actual load scale;
\item RL had no safe hybrid, no learning-disabled control and no sensitivity tests;
\item real OVS had never been executed.
\end{itemize}
Each was fixed or documented (\code{RUNTIME_PATH.md} \S3).
\end{qa}

\begin{qa}{802}{How do you know the earlier results were not an accident?}
Before any behavioural change, the frozen v4 evaluation was re-run on the same seeds under new run IDs: 240/240 runs.
Every per-condition difference was below the seed-to-seed spread, and the headline contrasts reproduced in size and
significance (\code{results/evidence/baseline_v4/BASELINE.json}).
\end{qa}

\section{Claims, tested once}

\begin{table}[!htbp]\centering\footnotesize
\begin{tabular}{@{}l p{4.6cm} r r@{\,}l r r l@{}}
\toprule
& claim (comparison) & $n$ & $\Delta$ & [95\,\% CI] & $r$ & $p_{\text{Holm}}$ & verdict\\
\midrule
C1 & critical class protected (qos\_adaptive vs proposed, healthcare SLA) & \NFCOneN & \NFCOneDiff & [\NFCOneLo, \NFCOneHi] & \NFCOneR & \NFCOnePHolm & \NFCOneVerdict\\
C2 & worst class not hurt (same, worst-class SLA; margin $-$0.01) & \NFCTwoN & \NFCTwoDiff & [\NFCTwoLo, \NFCTwoHi] & \NFCTwoR & -- & \NFCTwoVerdict\\
C3a & context helps the critical class (aware vs blind) & \NFCThreeaN & \NFCThreeaDiff & [\NFCThreeaLo, \NFCThreeaHi] & \NFCThreeaR & \NFCThreeaPHolm & \NFCThreeaVerdict\\
C3b & context helps after changes (transient SLA) & \NFCThreebN & \NFCThreebDiff & [\NFCThreebLo, \NFCThreebHi] & \NFCThreebR & \NFCThreebPHolm & \NFCThreebVerdict\\
C4 & context redundant when demand fits (equivalence $\pm$0.01) & \NFCFourN & \NFCFourDiff & [\NFCFourLo, \NFCFourHi] & \NFCFourR & -- & \NFCFourVerdict\\
C5 & reaction helps after changes & \NFCFiveN & \NFCFiveDiff & [\NFCFiveLo, \NFCFiveHi] & \NFCFiveR & \NFCFivePHolm & \NFCFiveVerdict\\
C6a & prediction helps after changes & \NFCSixaN & \NFCSixaDiff & [\NFCSixaLo, \NFCSixaHi] & \NFCSixaR & \NFCSixaPHolm & \NFCSixaVerdict\\
C6b & prediction helps overall & \NFCSixbN & \NFCSixbDiff & [\NFCSixbLo, \NFCSixbHi] & \NFCSixbR & \NFCSixbPHolm & \NFCSixbVerdict\\
C7 & residual RL beats the heuristic & \NFCSevenN & \NFCSevenDiff & [\NFCSevenLo, \NFCSevenHi] & \NFCSevenR & \NFCSevenPHolm & \NFCSevenVerdict\\
C8 & safe RL hybrid beats the heuristic & \NFCEightN & \NFCEightDiff & [\NFCEightLo, \NFCEightHi] & \NFCEightR & \NFCEightPHolm & \NFCEightVerdict\\
C9 & online learning beats learning-disabled & \NFCNineN & \NFCNineDiff & [\NFCNineLo, \NFCNineHi] & \NFCNineR & \NFCNinePHolm & \NFCNineVerdict$^{*}$\\
C10 & SliceNet everywhere non-inferior (margin $-$0.02) & \NFCOneZeroN & \NFCOneZeroDiff & [\NFCOneZeroLo, \NFCOneZeroHi] & \NFCOneZeroR & -- & \NFCOneZeroVerdict\\
C11 & self-healing cuts core outage (s) & \NFCOneOneN & \NFCOneOneDiff & [\NFCOneOneLo, \NFCOneOneHi] & \NFCOneOneR & \NFCOneOnePHolm & \NFCOneOneVerdict\\
\bottomrule
\end{tabular}
\caption{Predeclared claims (\capcode{results/tables/final_claims.csv}). $^{*}$See Q807 for a confound and its
replication.}
\label{tab:finalclaims}
\end{table}

\begin{qa}{803}{Do the requirements hold?}
\begin{itemize}
\item \textbf{R1, feasible SLAs: \NFROnePass.} All four classes reach $\ge$0.95 SLA share in every verified-feasible run
(\NFROneChecked{} runs, \NFROneFailures{} failures).
\honest{The plan listed school and harvest as ``feasible'', but school's downlink demand (130\,Mb/s) exceeds
100\,Mb/s. R1's own rule (count only runs whose every decision was SLA-feasible) excluded them, so R1 rests on the
normal-context runs.}
\item \textbf{R2, safety: \NFRTwoPass.} No failed verification, no planned capacity or protected-minimum violation and
no validator rejection in \NFRTwoChecked{} runs.
\item \textbf{R3, no starvation: \NFRThreePass} (\NFRThreeChecked{} contended runs of the QoS-targeted schedulers).
\end{itemize}
\end{qa}

\begin{qa}{804}{Why did C1 fail when the earlier evaluations showed a strong healthcare gain?}
Because of the everyday degradation policy introduced in this round, which was declared before the evaluation and
then not tuned after it. In everyday contexts it rations shortfalls \emph{proportionally}. On the 50\,Mb/s harvest
uplink, agriculture's 80\,Mb/s and healthcare's 5\,Mb/s compete, so healthcare loses the same fraction as
agriculture: its SLA share falls to \NFCHarvestFiveZeromAdHc{} (proposed \NFCHarvestFiveZeromPropHc).

Under the emergency policy (\emph{priority}) the gains are intact:
\begin{itemize}
\item disaster: \NFCDisasterPropHc{} $\rightarrow$ \NFCDisasterAdHc;
\item medical emergency: \NFCMedicalEmergencyPropHc{} $\rightarrow$ \NFCMedicalEmergencyAdHc;
\item night alert: \NFCMedicalNightPropHc{} $\rightarrow$ \NFCMedicalNightAdHc.
\end{itemize}
\honest{A protected minimum for the critical class under proportional rationing would likely fix harvest@50. It is an
algorithm change and therefore needs a new recorded evaluation. It was not made after seeing this result.}
\end{qa}

\begin{qa}{805}{Does context help?}
Where it carries information, yes, for the critical class. Context-aware vs context-blind raises healthcare SLA share
by \NFCThreeaDiff{} ($p_{\text{Holm}}=\NFCThreeaPHolm$) in the emergencies and timelines. It does not improve QoS in
the 10\,s after a change (C3b), and where demand fits it is redundant, as predicted (C4, equivalent within $\pm$0.01).
\trade{The blind controller rations proportionally everywhere, so agriculture fares better in emergencies (disaster
\NFCDisasterBlindAg{} vs \NFCDisasterAdAg) while healthcare fares worse. Context decides \emph{who} gets the scarce
capacity; it does not create capacity.}
\end{qa}

\begin{qa}{806}{Does prediction help?}
No. Reacting within a second (C5) and predicting ramps, queue growth and clock-driven context changes (C6) give no
measurable gain on gradual ramps, abrupt surges or context changes. Even the oracle, which reads the scripted timeline,
does not beat reactive control, so there is little left for a forecaster to win.
\begin{itemize}
\item Forecast quality: the reactive mean has the lowest error (MAE \NFqReactiveMae\,Mb/s vs Holt \NFqHoltMae{} and
ramp \NFqRampMae).
\item Ramp extrapolation lowers under-provisioning for surges but raises it for gradual ramps.
\end{itemize}
Prediction is \textbf{not} claimed as a contribution.
\end{qa}

\begin{qa}{807}{Does RL work?}
\begin{itemize}
\item \textbf{In simulation, yes.} The residual PPO beats its heuristic by +\NResDefaultRelMin--\NResDefaultRelMax\,\%
on held-out test seeds, for all three seeds. It holds under a heavier workload (1.3$\times$:
+\NResDefaultLoadMin--\NResDefaultLoadMax\,\%), but not on the held-out medical-emergency context
(\NResDefaultCtxMin{} to \NResDefaultCtxMax\,\%). Doubling the SLA weight raises the in-distribution gain but still
loses on the held-out context; quadrupling the loss weight behaves like the default.
\item \textbf{On the network, no.} \code{rl_residual} and the safe hybrid are statistically indistinguishable from the
heuristic (C7, C8).
\item \textbf{Online learning (C9) is supported against its learning-disabled control:} the same exploration without
updates costs \NFCNineDiff{} mean class SLA. But the predeclared run started from a checkpoint that already held
$\approx$500 samples per context from 40 earlier live runs (older code). A post-hoc replication from an empty
checkpoint gives \NPostFreshDiff{} [\NPostFreshLo, \NPostFreshHi] ($p=\NPostFreshP$, $n=\NPostFreshN$) against the same
control.
\end{itemize}
\honest{C9 compares learning with noisy exploration. It does not show that online RL beats the heuristic; that
comparison was not predeclared, and the run durations differ (60 vs 40\,s).}
\end{qa}

\section{Real Open vSwitch}

\begin{qa}{808}{Do the configured queues change real traffic outcomes?}
Yes, on real OVS bridges (privileged project image, Open vSwitch 3.3), seeds 30--32:
\begin{itemize}
\item disaster healthcare SLA share: \NOvsDisasterNoneHc{} without slicing, \NOvsDisasterStatHc{} static,
\NOvsDisasterAdHc{} QoS-targeted (agriculture \NOvsDisasterNoneAg{} $\rightarrow$ \NOvsDisasterAdAg);
\item school worst-class SLA: \NOvsSchoolNoneMin{} $\rightarrow$ \NOvsSchoolStatMin{} $\rightarrow$ \NOvsSchoolAdMin.
\end{itemize}
The OVS-native queue backend verified \NOvsBackendChanges{} of \NOvsBackendChanges{} rate changes by reading the OVS
queue database back (\NOvsBackendFailed{} failures).
\end{qa}

\begin{qa}{809}{Does self-healing work on real OVS?}
\begin{itemize}
\item Fiber cut: \NOvsFiberCutOutage\,s outage.
\item Repeated cuts: \NOvsRepeatedFiberCutOutage\,s, with \NOvsRepeatedFiberCutVerified{} verified recoveries per
run.
\item Backhaul degradation: switched before the cut, \NOvsBackhaulDegradationOutage\,s.
\item Core-switch failure: escalated, not healed (\NOvsCoreSwitchFailureOutage\,s).
\item No false failovers in \NOvsNoFaultRuns{} fault-free runs.
\end{itemize}
In emulation, self-healing cuts core-bound outage by \NFCOneOneDiff\,s against none (C11).
\honest{No radio-resource slicing and no live 5G RAN is claimed. The 5G study uses Open5GS with a simulated RAN
(UERANSIM), and the evaluation here is Mininet emulation plus real OVS switching.}
\end{qa}

\section{Verdict}

\begin{qa}{810}{Claim by claim, what can you say to the panel?}
\begin{description}
\item[Supported]
  \begin{itemize}
  \item Safe, verified actuation (R2).
  \item Starvation-free degradation (R3).
  \item All four classes meet SLAs in verified-feasible conditions (R1, narrow).
  \item Context protects the critical class where informative (C3a) and is redundant where demand fits (C4).
  \item Online learning beats learning-disabled exploration (C9, see Q807).
  \item SliceNet on every flow is non-inferior (C10).
  \item Self-healing (C11; real OVS).
  \end{itemize}
\item[Partially supported]
  \begin{itemize}
  \item QoS-targeted allocation protects the critical class in emergencies but not in everyday overload (C1).
  \item The worst-off class is not hurt within the stated margin only in some conditions (C2 inconclusive).
  \end{itemize}
\item[Unsupported]
  \begin{itemize}
  \item Context after changes (C3b).
  \item Reaction and prediction (C5, C6).
  \item Offline RL and the safe hybrid on the network (C7, C8).
  \end{itemize}
\item[Not evaluated] Radio-resource slicing; live 5G RAN.
\end{description}
\end{qa}
